\mathplain
\section{Introduction}\label{sec:intro}

Dependent Object Types (DOT) is the calculus behind Scala's type system \cite{amin2016essence}. A DOT object can have type members as well as term members, and a type can name the type member of a variable. If $x$ is an object with a member $A$, then $x.A$ is a type, and what is known about $x.A$ depends on what is known about $x$. That is the sense in which the types are dependent.

This report is about one theorem: well-typed DOT programs do not get stuck. The theorem has been proven several times \cite{amin2016essence,rompf2016soundness,rapoport2017simple,amin2017definitional,giarrusso2020scala}, and each proof had to deal with the same difficulty, called bad bounds. A type member $A$ is declared with a lower and an upper bound, $\dtyp{A}{S}{T}$, and nothing in DOT forces $S$ to be a subtype of $T$. Under a binder $x : \dtyp{A}{\top}{\bot}$ the type system derives $\top \sub x.A \sub \bot$, so every type is a subtype of every other. The textbook proof of type safety needs to invert subtyping, and in such a context there is nothing to invert.

The proof explained here takes a route that compilers take. GHC does not check its optimizations against the derivations of the Haskell type checker. It translates every program into System FC \cite{sulzmann2007fc}, an intermediate language in which every use of a type equality is an explicit proof term, a coercion, and a small checker for FC, Core Lint, can re-check the output of every pass. We do the same for DOT. A DOT typing derivation is translated into a term of a calculus called \FCdot, in which every use of subtyping is an explicit proof term. We call these proof terms evidence. Evidence erases to nothing, so the translated program computes what the source program computes. Type safety for \FCdot is then proven once, by preservation and progress, and type safety for DOT follows by running the source program and its translation side by side.

The translation is not the point. The point is that \FCdot has a metatheory with a shape the textbook reader will recognize. There is one lemma about evidence, and it replaces every inversion lemma about subtyping: over a typed store, closed evidence normalizes. A closed piece of evidence for an inclusion $S \le T$ normalizes to a head form, which says how $S$ is included in $T$: $S$ is $\bot$, or $T$ is $\top$, or both are function types and the form carries the coercions of the domain and the codomain, or both are object types and the form carries one entry per proposition of the object type $T$, and so on. The proof is a structural induction on the typing of the evidence. Bad bounds are not ruled out by the type system. Evidence that goes through bad bounds normalizes like all other evidence, and when it is closed over a store it cannot exist, because no head form fits it.

\paragraph{What this report covers.}
Section~\ref{sec:dot} introduces DOT for a reader who knows System F. Section~\ref{sec:problem} recalls the textbook proof of type safety and shows where bad bounds break it. Section~\ref{sec:miniature} is the core of the report. It develops the whole proof on a fragment with functions and type members only, small enough to state every definition and to work through examples by hand. The sections after it scale the fragment up. Section~\ref{sec:target} gives the full target \FCdot, Section~\ref{sec:canonical} the normalization theorem, and Section~\ref{sec:metatheory} preservation, progress, and erasure. Section~\ref{sec:translation} translates DOT into \FCdot, and Section~\ref{sec:safety} derives the safety of DOT, for a store machine and for the reduction semantics of the WadlerFest presentation. Section~\ref{sec:comparison} compares the proof with the direct proof of Rapoport et al.\ \cite{rapoport2017simple}. Section~\ref{sec:history} records the designs that did not work and why, Section~\ref{sec:mechanization} is a tour of the \Lean development, and Sections~\ref{sec:related} and~\ref{sec:conclusion} discuss related work and what is left out.

A reader who wants the idea and not the details can read Sections~\ref{sec:dot} to~\ref{sec:miniature} and then Section~\ref{sec:safety}. A reader who already knows DOT can skip Section~\ref{sec:dot} and skim Section~\ref{sec:problem}.

\paragraph{What is proven.}
Everything stated as a theorem or lemma in this report is mechanized in \Lean, with no library dependencies, and the main theorems depend on propositional extensionality and quotient soundness only. Section~\ref{sec:mechanization} maps each result to the declaration that proves it. The miniature of Section~\ref{sec:miniature} is a teaching device and is not a separate development. Its target is a fragment of \FCdot, and its results are the full results restricted to that fragment. The source calculus is the calculus of Amin et al.\ \cite{amin2016essence}, published in the Festschrift for Philip Wadler and called the WadlerFest calculus below, with variable paths only. Nothing else about that calculus is restricted: the body of a recursive type may be any type, and so may the operands of an intersection.

\paragraph{Conventions.}
Source syntax and judgments are set in \src{blue}, target syntax and judgments in \tgt{red}, and the syntax of the untyped runtime in black. Lean names are set in \lean{typewriter}. All syntax is scoped by de Bruijn indices in the mechanization. The report names every binder instead.
